\documentclass[11pt]{article}

\usepackage[margin=1in]{geometry}
\usepackage{setspace}
\usepackage{hyperref}

\setstretch{1.1}

\title{\textbf{C-SQD Go-to-Market Strategy}\
General Epistemic Audit Infrastructure}
\author{}
\date{}

\begin{document}

\maketitle

\section*{Executive Summary}

The central strategic question for C-SQD is not how to build a better journal, nor how to persuade researchers to abandon existing publication venues. The question is whether there exists demand for structured epistemic auditing: organized, transparent, evidence-based scrutiny of important scientific and technical claims.

The proposed go-to-market strategy is therefore to enter the market not as a publisher, but as infrastructure for decomposed epistemic audits. Publishing may eventually become an important application of the platform, but it should not be the initial wedge.

The initial objective is to become the place where organizations commission, conduct, and consume structured audits of claims, models, datasets, reports, and scientific results.

\section{The Wrong Starting Point}

A common instinct when building scientific review infrastructure is to compete directly with journals. This approach faces several difficulties.

Researchers are deeply embedded within existing publication ecosystems. Journals possess established prestige, institutional recognition, indexing, and incentive structures. Convincing authors to publish on a new platform requires overcoming substantial coordination problems before meaningful network effects can emerge.

More importantly, framing C-SQD as a publishing platform obscures its broader purpose. The underlying capability is not publication. The capability is structured scrutiny.

Scientific publication is merely one domain in which scrutiny is valuable.

\section{The Core Thesis}

The fundamental thesis is that important claims should be auditable.

Organizations routinely make decisions based on scientific and technical claims whose validity is uncertain. Existing mechanisms for evaluating such claims are fragmented, informal, expensive, and often opaque.

Examples include:

\begin{itemize}
\item biotechnology companies evaluating academic findings prior to licensing or collaboration;
\item venture capital firms evaluating technical startups;
\item foundations assessing scientific programs or grant proposals;
\item journals seeking external review;
\item universities evaluating controversial or high-profile work;
\item policymakers relying on scientific evidence for public decisions.
\end{itemize}

In each case, the underlying need is similar. An organization wishes to know whether a claim survives rigorous scrutiny.

C-SQD should therefore be framed as infrastructure for conducting and documenting that scrutiny.

\section{The Initial Customer}

The first customer is unlikely to be an individual researcher.

The more natural customer is an organization that already spends money on diligence, review, validation, or assessment.

These organizations have both the incentive and the budget to commission structured audits.

The value proposition is straightforward:

\begin{quote}
If a claim is important enough to influence decisions, it is important enough to audit.
\end{quote}

This framing is easier to communicate and easier to monetize than asking authors to migrate away from existing publication systems.

\section{The Initial Vertical}

The platform should begin within a domain where founder credibility and reviewer access already exist.

Computational biology, biomedical machine learning, and adjacent areas provide several advantages.

First, existing professional networks can be leveraged to recruit early reviewers.

Second, these communities are accustomed to engaging with quantitative evidence, benchmarking, reproducibility, and methodological criticism.

Third, many claims within these domains have substantial downstream consequences for biotechnology, healthcare, and investment decisions.

The goal is not immediate scale. The goal is a small number of high-quality audits conducted by respected participants.

\section{Commissioned Audits and Decomposed Judgment}

The primary market mechanism should not be framed as a traditional review bounty.

Unlike software vulnerabilities, scientific and technical evaluations rarely reduce to discrete findings that can be mechanically verified. Most important assessments require judgment, interpretation, and synthesis across multiple dimensions of evidence.

Accordingly, C-SQD should frame its core offering as commissioned epistemic audits.

A sponsor commissions an audit of a scoped claim, claim-warrant bundle, model, dataset, benchmark, grant proposal, technical report, or artifact-attached scientific claim. In the academic domain, a paper may be the entry point, source material, or discovery surface, but the audit target is the claim and the warrants by which attached artifacts are supposed to support, challenge, or limit it. Rather than soliciting a single overall review, the platform decomposes the evaluation into multiple independent assessment dimensions.

Examples might include:

\begin{itemize}
\item validity of causal identification;
\item statistical methodology;
\item benchmark design and evaluation;
\item reproducibility;
\item adequacy of controls;
\item fidelity of citations;
\item strength of evidential support;
\item relevance to the sponsor's decision context.
\end{itemize}

Each dimension is reviewed independently before broader conclusions are formed.

This design is inspired by the concept of mediating assessments developed in the decision-quality literature. The central insight is that complex judgments become more reliable when evaluators first assess intermediate questions separately rather than beginning with an overall verdict.

C-SQD operationalizes this principle through structured ElementReviews.

An ElementReview represents a focused assessment of a specific evaluative dimension. Reviewers provide evidence, reasoning, confidence estimates, limitations, and recommendations within a narrowly defined scope.

Only after the ElementReviews have been completed does the platform construct a broader SynthesisReview that integrates the individual assessments into an overall audit.

The resulting process transforms review from an exchange of global opinions into a structured process of decomposed judgment.

\section{Why Decomposition Matters}

The principal innovation of C-SQD is not simply that reviews are public or compensated. The principal innovation is the decomposition of evaluation itself.

Traditional peer review often encourages reviewers to form an overall impression early and then justify that impression through selective reasoning. Different reviewers may focus on entirely different concerns, making disagreement difficult to interpret.

A decomposed audit creates a more transparent structure.

Instead of asking whether a paper is ``good,'' or how much the literature appears to support a claim, reviewers assess specific warrant links and evaluation dimensions independently. Areas of agreement and disagreement become visible. Confidence can vary across dimensions. Strong evidence in one area does not automatically compensate for weakness in another.

This structure benefits both sponsors and reviewers.

Sponsors obtain a more interpretable assessment of risk and uncertainty. Reviewers contribute expertise within well-defined domains rather than being forced into broad global judgments. Subsequent readers can inspect the reasoning process rather than merely the conclusion.

Most importantly, the audit becomes cumulative. Individual ElementReviews remain useful even when conclusions change, new evidence emerges, or future synthesis efforts reinterpret the underlying claim.

\section{Compensation and Incentives}

Compensation should primarily support participation in the audit process rather than reward only the discovery of specific flaws.

Sponsors fund audits.

Reviewers receive compensation for producing qualified ElementReviews within defined scopes and standards.

Additional awards may be granted for unusually important findings, replication efforts, methodological corrections, or other contributions that materially improve the quality of the audit.

The purpose of compensation is not to create adversarial criticism. The purpose is to support rigorous scrutiny.

This distinction is important. C-SQD should be understood as infrastructure for funded evaluation rather than a marketplace for negative opinions.

\section{Subscription Infrastructure}

Once audits are occurring regularly, subscription products become substantially easier to justify.

Universities may subscribe for reviewer management, training, audit workflows, and institutional participation.

Companies may subscribe for diligence infrastructure, expert discovery, and audit commissioning.

Foundations may subscribe for grant evaluation and program assessment.

Revenue therefore emerges from two complementary sources.

The first is transaction revenue generated by commissioned audits.

The second is subscription revenue generated by organizations that require recurring access to the infrastructure.

\section{Publishing as a Later Expansion}

Publishing should not be the first product.

Publishing becomes attractive only after the platform has established itself as the locus of high-quality auditing activity.

At that stage, the platform possesses something more valuable than publication infrastructure: accumulated trust.

The narrative then shifts from:

\begin{quote}
Publish your work on C-SQD.
\end{quote}

to

\begin{quote}
The most thoroughly audited work already lives on C-SQD.
\end{quote}

This is a much stronger strategic position because publication emerges from demonstrated audit quality rather than from attempts to replicate existing journal structures.

\section{Minimum Viable Product}

The MVP should focus on validating the market for decomposed epistemic audits.

Four components are sufficient.

First, a claim page containing metadata, source materials, discussion, and audit history.

For Academic Peer Review, paper pages should remain searchable and useful: a user should be able to see every audit that uses a paper as evidence, challenges it, or audits a claim made within it. This is a navigation principle, not a change in audit object. Papers are evidence artifacts; scoped claims are the epistemic targets.

Second, an ElementReview system that supports structured mediating assessments across multiple evaluation dimensions.

Third, a SynthesisReview system that integrates completed ElementReviews into a coherent audit.

Fourth, a sponsored-audit mechanism through which organizations can commission evaluations and compensate reviewers.

Many capabilities envisioned for the broader C-SQD ecosystem can be deferred.

Publishing workflows, journals, reputation systems, synthesis networks, Observatory integrations, citation metrics, and other advanced features may eventually prove valuable. However, they are not necessary to test the central hypothesis.

The critical question is whether organizations value structured, decomposed audits sufficiently to fund them.

If they do, the remainder of the ecosystem can be built atop a validated foundation.

\section*{Conclusion}

The most important strategic insight is that C-SQD may not initially be a publishing platform at all.

Its natural first identity is infrastructure for decomposed epistemic audits.

Organizations already spend substantial resources evaluating scientific and technical claims. Existing approaches rely heavily on holistic expert judgment, fragmented diligence processes, and opaque evaluation mechanisms.

C-SQD offers an alternative: structured audits built from mediating assessments, transparent reasoning, and explicit synthesis.

If organizations value such audits sufficiently to fund them, then publishing, synthesis networks, reviewer ecosystems, and broader knowledge infrastructure can emerge on top of a validated foundation.

Accordingly, the first objective should be simple: determine whether there is a market for decomposed epistemic audits.

\end{document}
